\section{Μικρο-Μετρήσεις Επίδοσης Υποσυστημάτων}
\label{sec:eval_microbench}

\subsection{Μεθοδολογία Μικρο-Μετρήσεων και Υπολογιστικοί Τελεστές GA}
Για την ακριβή μέτρηση της υπολογιστικής καθυστέρησης των επιμέρους θεμελιωδών αλγορίθμων σε απομονωμένο περιβάλλον, χρησιμοποιήθηκε το πλαίσιο στατιστικών μετρήσεων \texttt{Criterion.rs} της Rust. Όλες οι μετρήσεις εκτελέστηκαν με 100 δείγματα μετά από κατάλληλη φάση προθέρμανσης (warmup) για την εξάλειψη παρεμβολών του λειτουργικού συστήματος.

Στο Σχήμα~\ref{fig:ga_operators_bench} παρουσιάζονται οι μετρήσεις των εξελικτικών τελεστών του Orchestrator. Στο άνω μέρος του Πίνακα~\ref{tab:ga_operators_gene_store} καταγράφονται οι αντίστοιχες μέσες τιμές καθυστέρησης.

\begin{figure}[htbp]
    \centering
    \includegraphics[width=0.92\textwidth]{figures/fig6_ga_operators_benchmark.png}
    \caption{Μικρο-μετρήσεις τελεστών Γενετικού Αλγορίθμου (Criterion.rs): (α) Δειγματοληψία αρχικού πληθυσμού, (β) επιλογή επιζώντων $(\mu + \lambda)$ και (γ) γκαουσιανή μετάλλαξη και περιορισμός (clamping) για 100 απογόνους.}
    \label{fig:ga_operators_bench}
\end{figure}

../../evaluation/figures/table3_ga_operators_and_gene_store.tex

Όπως φαίνεται στον Πίνακα~\ref{tab:ga_operators_gene_store}, η επιλογή $(\mu + \lambda)$ για πληθυσμό 100 ατόμων διαρκεί μόλις $1.70\ \mu\text{s}$, ενώ ακόμα και για $P=2.000$ άτομα ολοκληρώνεται σε $46.32\ \mu\text{s}$. Η γκαουσιανή μετάλλαξη 100 παιδιών διαρκεί $5.15\ \mu\text{s}$, επιβεβαιώνοντας ότι η διαχείριση της εξέλιξης σε Rust καταναλώνει αμελητέο ποσοστό του χρόνου κύκλου συγκριτικά με την αεροδυναμική προσομοίωση.

\subsection{Χωρική Μνήμη Gene Store και Διανυσματικός Υπολογισμός Αποστάσεων}
Η εσωτερική μνήμη γονιδίων (Gene Store) του Orchestrator επιτρέπει τον άμεσο εντοπισμό ακριβών επιτυχιών (Tier 1 hits) και την εκτέλεση τοπικών $k$-NN αναζητήσεων.

Στο Σχήμα~\ref{fig:gene_store_bench} και στο κάτω μέρος του Πίνακα~\ref{tab:ga_operators_gene_store} αναλύονται οι επιδόσεις της μνήμης.

\begin{figure}[htbp]
    \centering
    \includegraphics[width=0.92\textwidth]{figures/fig6_gene_store_benchmark.png}
    \caption{Μικρο-μετρήσεις χωρικής μνήμης Gene Store: (α) Διανυσματικός υπολογισμός 10D ευκλείδειας απόστασης (SIMD AVX2), (β) χρόνος ακριβούς επιτυχίας (hit) έναντι αστοχίας (miss), (γ) κλιμάκωση $k$-NN αναζήτησης και (δ) σάρωση εκκαθάρισης TTL.}
    \label{fig:gene_store_bench}
\end{figure}

Χάρη στη διανυσματική επιτάχυνση SIMD (AVX2), ο υπολογισμός της 10D απόστασης εκτελείται σε μόλις $4.65\text{ ns}$. Η ανάκτηση ακριβούς επιτυχίας (Exact Cache Hit) απαιτεί $173.35\text{ ns}$, ενώ η αναζήτηση $k=10$ πλησιέστερων γειτόνων σε αποθήκη 1.000 διαμορφώσεων ολοκληρώνεται σε $13.12\ \mu\text{s}$.

\subsection{Κωδικοποίηση Καμπυλών Hilbert Πολλαπλής Διερεύνησης}
Η μετατροπή των 10D σημείων σε ακέραιες συντεταγμένες και η παραγωγή των περιστραμμένων κλειδιών Hilbert αξιολογήθηκε για διαφορετικές διαστάσεις.

Στο Σχήμα~\ref{fig:hilbert_encoding_bench} και στον Πίνακα~\ref{tab:hilbert_encoding_bench} παρουσιάζονται οι χρόνοι εκτέλεσης των στοιχειωδών πράξεων Hilbert.

\begin{figure}[htbp]
    \centering
    \includegraphics[width=0.92\textwidth]{figures/fig6_hilbert_encoding_benchmark.png}
    \caption{Μικρο-μετρήσεις κωδικοποίησης καμπυλών Hilbert: (α) Μετασχηματισμός 10D συντεταγμένων σε δεκαεξαδικό συμβολοσειρά, (β) αποσειριοποίηση κλειδιού, και (γ) κλιμάκωση παραγωγής περιστραμμένων κλειδιών Multi-Probe ($C=3$) σε συνάρτηση με τις διαστάσεις του χώρου ($d=5, 10, 15$).}
    \label{fig:hilbert_encoding_bench}
\end{figure}

../../evaluation/figures/table5_hilbert_spatial_indexing.tex

Ο καθαρός μετασχηματισμός 10 διαστάσεων απαιτεί $505.2\text{ ns}$, ενώ η πλήρης παραγωγή 3 περιστραμμένων καμπυλών (multi-probe key generation) για $d=10$ διαρκεί $2.18\ \mu\text{s}$, επιτρέποντας στον κόμβο LEAD να αναζητά χωρικούς γείτονες χωρίς να επιβαρύνει τη ροή του εξελικτικού αλγορίθμου.

\subsection{Επανεκπαίδευση RMI και Ομόσπονδος Συγχρονισμός}
Η συντήρηση του ευρετηρίου RMI σε περιβάλλον συνεχούς εισροής δεδομένων περιλαμβάνει επιγραμμική προσαρμογή των φυλλο-προβλεπτών, αυτόματο συντονισμό (auto-tuning) και ομόσπονδο συγχρονισμό (Federated Averaging) μεταξύ των κόμβων LEAD DHT.

Στο Σχήμα~\ref{fig:rmi_train_bench} και στον Πίνακα~\ref{tab:rmi_training_sync} παρουσιάζονται οι μετρήσεις εκτέλεσης αυτών των διεργασιών παρασκηνίου.

\begin{figure}[htbp]
    \centering
    \includegraphics[width=0.92\textwidth]{figures/fig6_rmi_train_benchmark.png}
    \caption{Μικρο-μετρήσεις επανεκπαίδευσης RMI και ομόσπονδου συγχρονισμού: (α) Γραμμική προσαρμογή φύλλων RMI σε συνάρτηση με το πλήθος κλειδιών, (β) αλγόριθμος ορειβατικού συντονισμού (Mountain-Climbing Auto-Tune), (γ) ομόσπονδη συνένωση βαρών FedAvg για $M=3, 8, 16$ κόμβους DHT, και (δ) 2-bit βρόχος προσαρμογής PID.}
    \label{fig:rmi_train_bench}
\end{figure}

../../evaluation/figures/table4_rmi_training_and_federated_sync.tex

Η προσαρμογή των γραμμικών προβλεπτών RMI για 500 κλειδιά εκτελείται σε μόλις $53.5\ \mu\text{s}$, επιτρέποντας την εκτέλεσή της εν ριπή οφθαλμού χωρίς διακοπή της υπηρεσίας ανάγνωσης. Επιπρόσθετα, η ομόσπονδη συνένωση παραμέτρων (FedAvg) για τους 3 κόμβους του δακτυλίου ολοκληρώνεται σε $210.4\text{ ns}$, ενώ ο βρόχος προσαρμογής PID (2-bit anchor adjustment) απαιτεί μόλις $5.52\text{ ns}$, καθιστώντας δυνατή την αποδοτική διατήρηση του ευρετηρίου ακόμη και υπό υψηλούς ρυθμούς εγγραφής νέων διαμορφώσεων.
